\subsection{\texorpdfstring{空间 \(L_{F}^{1}\) 的对偶（F 为可分 Banach 空间）}{空间 L sub F superscript 1 的对偶（F 为可分 Banach 空间）}}

命题 10. --- 设 F 是可分 Banach 空间。对每个函数 \(\mathbf{f} \in \overline{\mathcal{L}}_{\mathrm{F}}^{1}\) 与每个函数 \(\mathbf{g} \in \mathcal{L}_{\mathrm{F}'_s}^{\infty}\) ，数值函数 \(\langle \mathbf{f}, \mathbf{g} \rangle : t \mapsto \langle \mathbf{f}(t), \mathbf{g}(t) \rangle\) 本质 \(\mu\)-可积，且

\[
\left| \int \langle \mathbf {f}, \mathbf {g} \rangle d \mu \right| \leqslant \overline {{{{N}}}} _ {1} (\mathbf {f}) N _ {\infty} (\mathbf {g}).\tag{3}
\]

对每个类 \(\dot{\mathbf{g}}\in\mathrm{L}_{\mathrm{F}^{\prime}s}^{\infty}\) ，令 \(\theta(\dot{\mathbf{g}})\) 为 \(\mathrm{L}_{\mathrm{F}}^{1}\) 上的连续线性形式，它是由线性形式 \(\mathbf{f}\mapsto\int\langle\mathbf{f},\mathbf{g}\rangle d\mu\) （ \(\overline{\mathcal{L}}_{\mathrm{F}}^{1}\) 上的）经取商得到的；这时 \(\theta\) 是 \(\mathrm{L}_{\mathrm{F}^{\prime}s}^{\infty}\) 到强对偶 \((\mathrm{L}_{\mathrm{F}}^{1})^{\prime}\) （Banach 空间 \(\mathrm{L}_{\mathrm{F}}^{1}\) 的）上的线性等距。

对 T 的每个紧子集 K 与每个 \(\varepsilon > 0\) ，存在 K 的紧子集 \(K'\) ，使得 \(\mu(K - K') \leqslant \varepsilon\) ，且 f（相应地， g）到 \(K'\) 的限制是 \(K'\) 到 F（相应地，到 \(F'_s\) ）中的连续映射；因此 \(\mathbf{g}(K')\) 是弱紧的，从而在 \(F'\) 上等度连续（TVS, III, §4, No. 2, 定理 1 或 IV, §1, 习题 10）。而今， \(F \times F'\) 上的典范双线性形式到 F 与 \(F'\) 的等度连续子集的乘积上的限制，对 F 的拓扑与 \(\sigma(F', F)\) 的乘积拓扑是连续的（GT, X, §2, No. 1, 命题 1 的推论 4）；由此， \(\langle f, g \rangle\) 到 \(K'\) 的限制连续，从而 \(\langle f, g \rangle\) 是 \(\mu\)-可测的。此外，

\[
| \langle \mathbf {f} (t), \mathbf {g} (t) \rangle | \leqslant | \mathbf {f} (t) | \cdot | \mathbf {g} (t) | \leqslant | \mathbf {f} (t) | N _ {\infty} (\mathbf {g})
\]

局部几乎处处成立，从而 \(\langle f,g\rangle\) 本质 \(\mu\)-可积且不等式 (3) 成立（Ch. IV, §5, No. 6, 定理 5 及 Ch. V, §1, No. 3, 引理）。

剩下要证 \(\theta\) 是满射等距。设 u 是 \(L_{F}^{1}\) 上的连续线性形式。映射 \((h,\mathbf{z})\mapsto u(\widetilde{h}\mathbf{z})\) 是 \(L^{1}\times F\) 上的连续双线性形式，因为

\[
\left| u \left(\widetilde {h} \mathbf {z}\right) \right| \leqslant \| u \| \cdot \mathrm{N} _ {1} (h \mathbf {z}) \leqslant \| u \| \cdot | \mathbf {z} | \cdot \mathrm{N} _ {1} (h).
\]

由第 5 目定理 1 的推论 1，存在 T 到 F' 中的映射 g ，属于 \(L_{F_{s}^{\infty}}\) ，使得 \(|\mathbf{g}(t)| \leqslant \|u\|\) （对所有 \(t \in T\) ），且 \(u(\widetilde{h}\mathbf{z}) = \int \langle h\mathbf{z}, \mathbf{g} \rangle d\mu\) （对每个函数 \(h \in L^{1}\) ——其类 \(\widetilde{h}\) 在 \(L^{1}\) 中——与每个 \(z \in F\) ）。换言之，线性形式 u 与 \(\theta(\dot{\mathbf{g}})\) 在 \(L_{F}^{1}\) 的由形如 \(\widetilde{h}\mathbf{z}\) 的元素（ \(\widetilde{h} \in L^{1}\) , \(z \in F\) ）生成的子空间上一致。由于该子空间在 \(L_{F}^{1}\) 中稠密（Ch. IV, §3, No. 5, 命题 10），得出 \(u = \theta(\dot{\mathbf{g}})\) ，这已经证明 \(\theta\) 是满射的。此外，由 (3)， \(\|\theta(\dot{\mathbf{g}})\| \leqslant N_{\infty}(\mathbf{g}) \leqslant \|u\| = \|\theta(\dot{\mathbf{g}})\|\) ，从而 \(\|\theta(\dot{\mathbf{g}})\| = N_{\infty}(\mathbf{g})\) ，证明完毕。
% semantic-fragments: legacy-input-seam
\relax{}
